\Askhsh{1794}{4}
\begin{ylh}{1}
    \item 5.3. Ορθογώνιο
    \item 5.4. Ρόμβος
    \item 5.6. Εφαρμογές στα τρίγωνα.
\end{ylh}
\begin{alist}
    \item Σε ορθογώνιο $\mathrm{A}\mathrm{B}\Gamma\Delta$ θεωρούμε $\mathrm{K}, \Lambda, \mathrm{M}, \mathrm{N}$ τα μέσα των πλευρών του $\mathrm{A}\mathrm{B}, \mathrm{B}\Gamma, \Gamma\Delta, \Delta\mathrm{A}$ αντίστοιχα. Να αποδείξετε ότι το τετράπλευρο $\mathrm{K}\Lambda\mathrm{M}\mathrm{N}$ είναι ρόμβος. \monades{15}
    \item Σε ένα τετράπλευρο $\mathrm{A}\mathrm{B}\Gamma\Delta$ τα μέσα $\mathrm{K}, \Lambda, \mathrm{M}, \mathrm{N}$ των πλευρών του $\mathrm{A}\mathrm{B}, \mathrm{B}\Gamma, \Gamma\Delta, \Delta\mathrm{A}$ αντίστοιχα είναι κορυφές ρόμβου. Το τετράπλευρο $\mathrm{A}\mathrm{B}\Gamma\Delta$, πρέπει να είναι απαραίτητα ορθογώνιο; Να τεκμηριώσετε τη θετική ή αρνητική σας απάντηση. \monades{10}
\end{alist}
